\chapter{\kesimpulan}
\label{bab:6}
This chapter will discuss the summary of this research in building a GraphRAG framework consisting of only open-source components. Additionally, several potential future works are also addressed to improve this research in the future.

\section{Conclusion}
In this research, we implement a system for performing RAG over KGs using only open-source components. We design and develop two pipelines. The first pipeline is simpler but has a narrower scope, focusing exclusively on Wikidata. It consists of several stages: entity extraction and linking, SPARQL query generation (using the top-100 most frequently used properties augmented into the prompt), and answer generation. However, this approach limits the types of questions that the system can handle, since the properties are only obtained from the pre-defined top-100 properties list. To overcome this limitation, we introduce a second pipeline that builds upon the first. It incorporates verbalization-based retrieval to handle single-hop questions and integrates a vector database for flexible retrieval of ontological elements (properties and classes) without requiring explicit mention of all elements in the prompt. This enhancement significantly improves the system's ability to support a broader range of question types.

The final experiment demonstrates that the usage of Qwen2.5 7B LLM in the second pipeline serves as the best configuration for all datasets: Wikidata, DBpedia, and the local curriculum KG. This configuration yields accuracies of $0.458$, $0.517$, and $0.805$ for Wikidata, DBpedia, and the curriculum KG, respectively. Our findings indicate that simpler KGs (and ontology) tend to yield better performance, as observed with the curriculum KG. Moreover, the multilingual support experiment demonstrated that the system could process queries in Indonesian, achieving an accuracy of $0.362$. While this score is lower than that of the English dataset, it highlights the system's foundational ability to handle multilingual queries through translation and retrieval mechanisms.

Although the results for Wikidata and DBpedia appear relatively low, we believe their potential accuracy could be higher. This conclusion stems from our observation of several inaccuracies within these datasets, which likely impacted the evaluation outcomes. Furthermore, we found that the ablation study underscores the critical role of entity and property retrieval in providing context to the LLM, significantly impacting the system's overall performance. Removing this component caused major accuracy drops across all datasets, with the curriculum KG dataset being the most affected, where accuracy fell to zero. Another interesting finding is that for simpler KGs, the use of verbalization deteriorates pipeline performance. This decline is attributed to the misinterpretation of questions and limitations in handling aggregate and other complex operations.

Finally, the objectives from Subsection \ref{sec:tujuan} are completed throughout this research by:
\begin{enumerate}
    \item Developing a RAG system using KG as the primary knowledge base, which pipelines are explained in Subsection \ref{exp-detail} and Chapter \ref{bab:4}.
    \item Determining that Qwen2.5 7B is the best-performing LLM, as explained in Subsubsubsection \ref{ga-llm}.
    \item Identifying that the most influential component within the architecture is the entity and property retrieval, which information is obtained through the ablation study (Subsection \ref{ablation}).
\end{enumerate}

\section{Future Works}
Given the limitations of the current research, there are several possible future works.
\begin{enumerate}
    % \item Implementing the pipeline for the Bahasa Indonesia version. This would include developing a component capable of translating user questions into Bahasa Indonesia, ensuring the system can handle queries in the local language effectively. Alternatively, leveraging an LLM that natively supports Bahasa Indonesia could streamline the process and enhance performance.
    \item Fine-tuning the LLM for SPARQL query generation task. This would allow the model to better understand the nuances of query structure, entity relationships, and contextual mapping from natural language to SPARQL, ultimately improving accuracy and performance in generating precise queries.
    \item Exploring the use of LLMs with a larger number of parameters. Models with more parameters typically have enhanced capacity to understand and generate complex queries, potentially leading to improved performance in SPARQL query generation. However, this approach would also require addressing challenges related to computational efficiency and resource requirements.
    \item Introducing a dedicated classifier to determine whether a question can be resolved using verbalization (can be in the form of an intent classification task). This classifier would analyze the input question and decide whether verbalization is an appropriate approach to answer the question. By integrating this component, the system could optimize its query generation process, ensuring that verbalization is only applied when suitable, while alternative strategies are employed for more complex questions.
    \item Scaling the system to larger knowledge graphs. The current implementation is limited by memory constraints due to the use of \texttt{rdflib} for in-memory graph storage. Future work could involve optimizing the system to handle larger knowledge graphs by implementing a persistent storage solution, such as a database-backed storage approach (e.g., using SQLite, GraphDB, or Virtuoso), to efficiently manage larger datasets and improve scalability.
\end{enumerate}